\par\Needspace{4\baselineskip}
\begin{center}
\hypertarget{native-057-article039}{}\pdfbookmark[2]{39.}{article-039}39.
\end{center}
Offenbar ließe sich fast die Hälfte des Aufwands für dieses Probieren einsparen, wenn man die Tafel so einrichtete, dass sich $\log B$ unmittelbar mit dem Argument $A$ entnehmen ließe. Dann blieben drei Rechenschritte. Der erste ist indirekt und besteht darin, $A$ so zu bestimmen, dass Gleichung 1) aus Art. 37 erfüllt ist. Der zweite, die Bestimmung von $E$ aus $A$ und $B$, erfolgt direkt entweder mit der Gleichung $E=2B(A^{\frac{1}{2}}+\tfrac{1}{15}A^{\frac{3}{2}})$ oder mit der Gleichung $\sin E=2B(A^{\frac{1}{2}}-\tfrac{3}{5}A^{\frac{3}{2}})$. Der dritte ist die Bestimmung von $v$ aus $E$ nach Gleichung VII aus Art. 8. Den ersten Schritt führen wir auf einen raschen Algorithmus zurück, der ohne planloses Probieren auskommt. Den zweiten und dritten fassen wir zu einem einzigen zusammen, indem wir die neue Größe $C$ in unsere Tafel aufnehmen. So kommen wir ganz ohne $E$ aus und gewinnen zugleich eine elegante und zweckmäßige Formel für den Radiusvektor. Wir gehen diese Punkte nun der Reihe nach durch.

Zunächst formen wir Gleichung 1) so um, dass sich die \textsc{Barker}sche Tafel zu ihrer Lösung verwenden lässt. Dazu setzen wir $A^{\frac{1}{2}}=\operatorname{tang}\tfrac{1}{2}w\,\mathord{.}\,\BareRoot{\frac{5-5e}{1+9e}}$ und erhalten
\[
75\operatorname{tang}\tfrac{1}{2}w+25\operatorname{tang}\tfrac{1}{2}w^3=\frac{75kt\BareRoot{(\tfrac{1}{5}+\tfrac{9}{5}e)}}{2Bq^{\frac{3}{2}}}=\frac{\alpha t}{B},
\]
wobei wir die Konstante $\frac{75k\BareRoot{(\tfrac{1}{5}+\tfrac{9}{5}e)}}{2q^{\frac{3}{2}}}$ mit $\alpha$ bezeichnen\HistoricalNoteEight{Handschriftliche Bemerkung:] Wenn man \textsc{Burckhardts} Tafel [Connaissance des Temps 1818, siehe auch Bd. VI, S. 596] verwendet, gilt $\frac{\BareRoot{(\tfrac{1}{10}+\tfrac{9}{10}e)}}{q^{\frac{3}{2}}}=\alpha$.}. Wäre also $B$ bekannt, so ließe sich $w$ sogleich der \SourceLeaf{058}\textsc{Barker}schen Tafel entnehmen. Dort ist es die wahre Anomalie, die zur mittleren Bewegung $\frac{\alpha t}{B}$ gehört. Aus $w$ ergibt sich $A$ nach der Formel $A=\beta\operatorname{tang}\tfrac{1}{2}w^2$, wobei wir die Konstante $\frac{5-5e}{1+9e}$ mit $\beta$ bezeichnen. Zwar ergibt sich $B$ erst aus $A$ mithilfe unserer Hilfstafel. Da seine Abweichung von der Einheit aber sehr klein ist, werden $w$ und $A$ voraussichtlich nur mit einem geringen Fehler behaftet sein, wenn man den Divisor $B$ zunächst ganz außer Acht lässt. Wir bestimmen daher zuerst $w$ und $A$ überschlägig, indem wir $B=1$ setzen. Mit dem Näherungswert von $A$ entnehmen wir unserer Hilfstafel $B$ selbst und wiederholen damit dieselbe Rechnung genauer. Meist ergibt sich aus dem so korrigierten Wert von $A$ genau derselbe Wert von $B$ wie aus dem Näherungswert. Eine weitere Wiederholung der Rechnung ist dann überflüssig, außer wenn der Wert von $E$ bereits beträchtlich ist. Es braucht kaum erwähnt zu werden, dass man einen anderweitig bereits zu Beginn bekannten Näherungswert von $B$ am besten gleich bei der ersten Näherung verwendet. Das ist stets möglich, wenn von mehreren nahe beieinanderliegenden zu berechnenden Orten schon einer oder zwei berechnet sind. Wer geübt rechnet, braucht die Rechnung dann sehr oft kein einziges Mal zu wiederholen. Diese sehr rasche Näherung verdanken wir dem Umstand, dass $B$ von $1$ nur um eine Größe vierter Ordnung abweicht, die zudem mit einem sehr kleinen Zahlenkoeffizienten multipliziert ist. Nun wird deutlich, dass wir diesen Vorteil dadurch gewonnen haben, dass wir die Größen $E-\sin E$ und $\tfrac{9}{10}E+\tfrac{1}{10}\sin E$ anstelle von $E$ und $\sin E$ eingeführt haben.
